\documentclass{article}

% NeurIPS 2026 style file, unmodified, as required by the ML4PS 2026 guidelines.
% Submission mode (no [final]/[preprint] option) => anonymous, line-numbered.
\usepackage{neurips_2026}

% ML4PS 2026 requires the conference footer to be replaced with the workshop
% notice. This is the only permitted deviation; the template itself is untouched.
\makeatletter
\renewcommand{\@noticestring}{%
  Submitted to the 9th Workshop on Machine Learning and the Physical Sciences
  (ML4PS 2026). Do not distribute.%
}
\makeatother

\usepackage[T1]{fontenc}
\usepackage{hyperref}
\usepackage{url}
\usepackage{booktabs}
\usepackage{amsfonts}
\usepackage{amsmath}
\usepackage{graphicx}

\title{Where Classical Interpolation and a Convolutional Surrogate Fail:
Wavefunction Super-Resolution on the 1D Quantum Harmonic Oscillator}

\author{%
  Anonymous Author(s) \\
  Affiliation \\
  Address \\
  \texttt{email} \\
}

\begin{document}

\maketitle

\begin{abstract}
We train a 1D convolutional network to reconstruct 1024-point
ground-state wavefunctions of the quantum harmonic oscillator from 64-point
coarse solutions, and compare it with cubic-spline interpolation on 200 matched
held-out samples. Both clear the $F > 0.99$ fidelity threshold on every sample
(network mean 0.999832, spline 0.999790), so the comparison turns on how each
method's error is organized over the frequency range. The difference is sharpest
where the physics is easiest: in the lowest-frequency bin the spline attains its
lowest error while the network attains its highest, on identical inputs. Spline infidelity is perfectly
rank-correlated with frequency (Spearman $\rho = +1.000$), a family-agnostic
failure, whereas network infidelity follows an inverted-U over the sampled
interval, worst at both edges ($\rho = +0.861$ against distance from the
midpoint). Retraining on a wider interval moved error off the old edges but not
symmetrically onto the new ones, only partly supporting a training-density
explanation.
\end{abstract}

\section{Introduction}

Direct solution of the time-independent Schr\"odinger equation (TISE) on a fine
grid scales badly: an $N$-point 1D problem requires diagonalizing an $N \times N$
Hamiltonian, and dimension and particle number compound the cost exponentially. A
natural mitigation is to solve cheaply on a coarse grid and reconstruct the fine
solution with a learned model. We train a convolutional network, QuantumResNet,
to map a 64-point coarse ground state of the quantum harmonic oscillator to its
1024-point counterpart, scored by quantum fidelity.

The 1D oscillator does not need machine learning; it solves in tens of
milliseconds and offers exact ground truth across a densely sampled family. Since
a cubic spline clears $F > 0.99$ across the range, a higher mean fidelity would
say little; the informative question is what organizes each method's error: the
spline's by resolution, the network's by position within its training range. The model is a supervised
surrogate, not a physics-informed neural network: its loss is pure mean-squared
error against precomputed eigenvectors, with no Schr\"odinger-residual term, the
defining feature of a PINN \citep{raissi2019}. ``Quantum'' here means quantum
simulation on classical hardware, with no qubits or circuits.

\section{Related Work}

PINNs embed a PDE residual in the training loss \citep{raissi2019}; ours does
not. Neural operators such as the FNO learn resolution-independent maps with
zero-shot super-resolution \citep{li2021fno}; ours is resolution-specific but
reconstructs fine fields from coarse observations of a family indexed by
$\omega$, as in SRCNN \citep{dong2016}. Closest are convolutional surrogates for
the Schr\"odinger equation: \citet{mills2017} map a 2D potential to a
ground-state energy, whereas we reconstruct the eigenvector; \citet{li2025} apply
super-resolution-inspired upsampling to electron densities, scoring by energy
error. Neural-network quantum states parameterize the wavefunction variationally
\citep{carleo2017}. That learned surrogates degrade away from their training
distribution is well documented: \citet{mouli2024} report neural operators
failing on out-of-domain PDE inputs at high in-domain accuracy, and
\citet{gopakumar2026} give conformal coverage across a surrogate's operating
domain under an exchangeability assumption. What this setting adds is
separability: a classical baseline on identical inputs has error ordered by
physics without exception ($\rho = +1.000$, zero inversions across 200 samples),
so both failure modes are observable side by side.

\section{Methods}

\textbf{System and ground truth.} We solve the 1D TISE in atomic units
($\hbar = m = 1$), $-\tfrac{1}{2}\psi'' + V(x)\psi = E\psi$, with
$V(x) = \tfrac{1}{2}\omega^{2}x^{2}$ on $[-10, 10]$ with Dirichlet boundaries. A
three-point finite-difference stencil yields a sparse Hamiltonian; the ground
state comes from \texttt{eigsh(H, k=1, which="SA")} at $N = 1024$ (target) and
$N = 64$ (input). Wavefunctions are $L^{2}$-normalized with fixed sign; only the
nodeless $n = 0$ state is used, at fixed centre $x_{0} = 0$. The state is a Gaussian of width $\sigma = 1/\sqrt{\omega}$
against spacing $\Delta x = 20/63 \approx 0.317$, so coarse points per width fall
from 4.45 at $\omega = 0.5$ to 1.41 at $\omega = 5.0$.

\textbf{Data and training.} 2{,}000 samples, $\omega \sim U[0.5, 5.0]$, seed 42;
split 72/18/10 (1{,}440/360/200). QuantumResNet maps $[B, 1, 64] \to [B, 1,
1024]$ through four transposed-convolution stages (kernel 4, stride 2, padding
1), each followed by a pre-activation residual block (BatchNorm/GELU/Conv1d,
twice, with skip); 64 channels, 166{,}209 parameters. Adam
(lr $10^{-3}$, batch 64), MSE loss, cosine annealing over the 30-epoch budget;
early stopping (patience 7) restored the best-validation checkpoint. A run takes
113\,s on an Intel Core Ultra 7 256V (CPU). Early stopping did not fire,
which is not headroom: the schedule drives the learning rate to near zero by
epoch 30. Robustness checks retrain the same configuration under
seeds 43--46; one single-seed ablation uses $\omega \in [0.3, 7.0]$,
$N = 3{,}000$ (448 samples per unit $\omega$ vs the canonical 444).

\textbf{Metrics and baselines.} Fidelity is
$F = (\sum_{i} \psi_{\mathrm{hr},i}\psi_{\mathrm{pred},i}\Delta x)^{2}$, both
renormalized under the released code's discrete-sum rule; we report infidelity
$1 - F$. Baselines: a cubic spline
(\texttt{CubicSpline}, $64 \to 1024$); Fourier zero-padding applied to the 63
unique periodic points of the endpoint-inclusive grid; and two oracles that
assume the harmonic family. The \emph{moment oracle} estimates $\omega$ from the
coarse input through $\mathrm{var}(x) = 1/(2\omega)$ and emits the analytic state
$(\omega/\pi)^{1/4}\exp(-\omega x^{2}/2)$. The \emph{inversion oracle} recovers
$\omega$ by minimizing $1 - F$ between the input and the coarse
finite-difference state at $\omega$ (bounded Brent, \texttt{xatol} $10^{-8}$,
seeded from the moment estimate), then emits the fine-grid state there;
costing a fine solve per sample, it is not a candidate method but a bound on the
information the coarse input carries. None is tuned. Error shapes use
Spearman correlations of infidelity against $\omega$ and against
$|\omega - \omega_{\mathrm{mid}}|$ ($\omega_{\mathrm{mid}} = 2.750$); the second
predictor is a folded transform of the first, so the pair restates
Table~\ref{tab:bins} rather than testing between explanations.

\section{Results}
\label{sec:results}

\textbf{Fidelity.} On the 200-sample test set the network reached mean fidelity
0.999832 (median 0.999889, min 0.999328, std $1.13 \times 10^{-4}$), 0/200 below
$F = 0.99$; the spline reached 0.999790 (min 0.999442, std
$1.65 \times 10^{-4}$). Across seeds 42--46 network mean fidelity is 0.999869
(0.999832--0.999914); seed 42, the headline run, is the lowest.
The Fourier baseline initially scored 0.969008 through a grid-convention error,
not a property of Fourier interpolation: an endpoint-inclusive array transformed
as periodic in index space, shifting it 0.149 units. Corrected it scores
0.999770. The moment oracle
reaches 0.999789 and tracks the spline bin by bin rather than the network, its
$\omega$ estimate carrying a mean $|\Delta\omega|$ of 0.13. The inversion oracle
recovers $\omega$ to a mean $|\Delta\omega|$ of $6.09 \times 10^{-8}$ (max
$2.81 \times 10^{-7}$) and matches the target to a mean
infidelity below $1.1 \times 10^{-16}$ in every bin, at double-precision
round-off. The coarse input therefore determines the state completely within the
family, so the spline's losses at high $\omega$ reflect family-agnosticism and
the moment oracle's a mis-specified estimator, the analytic identity applied to a
finite-difference state, rather than missing information; the network's
$1.68 \times 10^{-4}$ lies far above what the input permits.

\begin{table}[t]
  \centering
  \caption{Fidelity and infidelity by $\omega$ bin; all methods on the same 200
  matched samples (seed 42). FFT is the corrected Fourier baseline.}
  \label{tab:bins}
  \scriptsize
  \setlength{\tabcolsep}{4pt}
  \begin{tabular}{lrcccccc}
    \toprule
    $\omega$ range & $n$ & Spline $F$ & FFT $F$ & Moment $F$ & CNN $F$ &
    Spline $1-F$ & CNN $1-F$ \\
    \midrule
    $[0.5, 1.4)$ & 36 & 0.999980 & 0.999979 & 0.999981 & 0.999680 & 2.03e$-$5 & 3.20e$-$4 \\
    $[1.4, 2.3)$ & 36 & 0.999932 & 0.999930 & 0.999936 & 0.999894 & 6.75e$-$5 & 1.06e$-$4 \\
    $[2.3, 3.2)$ & 46 & 0.999836 & 0.999826 & 0.999840 & 0.999896 & 1.64e$-$4 & 1.04e$-$4 \\
    $[3.2, 4.1)$ & 42 & 0.999708 & 0.999682 & 0.999709 & 0.999892 & 2.92e$-$4 & 1.08e$-$4 \\
    $[4.1, 5.0]$ & 40 & 0.999527 & 0.999466 & 0.999511 & 0.999777 & 4.73e$-$4 & 2.23e$-$4 \\
    \midrule
    Overall & 200 & 0.999790 & 0.999770 & 0.999789 & 0.999832 & 2.10e$-$4 & 1.68e$-$4 \\
    \bottomrule
  \end{tabular}
\end{table}

\textbf{Matched comparison.} The network wins 132 of 200 matched samples (66\%;
one-sided binomial $p = 3.5 \times 10^{-6}$; Wilcoxon on paired infidelities
$W = 6588$, $p = 2.4 \times 10^{-5}$), with $1.25\times$ lower mean infidelity.
Both clear $F > 0.99$ on every sample, so these are differences in the structure
of very small errors, not a pass/fail separation. The aggregate conceals a
reversal (Figure~\ref{fig:err}): the network loses all 36 samples in the lowest
bin and 32 of 36 in the second, then wins all 128 above; below $\omega = 1.5$ it
wins 0 of 42, and at $\omega \geq 3.5$ 65 of 65 with $2.31\times$ lower mean
infidelity. For this seed the spline's last win falls at $\omega = 2.03$ and the
network wins every sample from $\omega = 2.09$ up; across five seeds that run
begins between $\omega = 1.73$ and 2.10 (mean 1.95) with 122--157 wins of 200, so
the crossover is a band, not a fixed value.

\begin{figure}[t]
  \centering
  \includegraphics[width=0.76\linewidth]{figure1.png}
  \caption{(a) Infidelity against $\omega$ for spline and network on the 200
  matched samples (open markers) with bin means, log ordinate; the curves cross
  near $\omega \approx 2$ (shaded). (b) Per-sample ratio of spline to network
  infidelity, crossing unity at the same transition.}
  \label{fig:err}
\end{figure}

\textbf{Two error structures.} The clearest evidence is where the physics is
easiest. In the lowest bin the spline attains $2.03 \times 10^{-5}$, the lowest
error of any method in any bin, while the network attains
$3.20 \times 10^{-4}$, its own worst and three times its interior error. All three non-learned baselines attain their lowest
error in that bin, where the network attains its highest, so undersampling cannot
be what the network is failing on. Numerically, spline infidelity is perfectly
rank-correlated with $\omega$ ($\rho = +1.000$, zero inversions) and uncorrelated
with the folded predictor ($\rho = +0.031$, $p = 0.66$); network infidelity is
uncorrelated with $\omega$ ($\rho = +0.026$, $p = 0.71$) and strongly correlated
with distance from the midpoint ($\rho = +0.861$, $p = 4.0 \times 10^{-60}$).
Re-expressing the predictor as $\sigma = 1/\sqrt{\omega}$ or $\log\omega$ leaves
these unchanged up to sign, as monotone reparametrization must, excluding a
parametrization artefact. Low-edge infidelity exceeds
high-edge by $1.433\times$.

A distributional reading fits this: accuracy set by local training density, worst
where a sample has neighbours on one side only. Retraining once on
$\omega \in [0.3, 7.0]$ at matched density held that prediction only in part.
Both former edges improved once interior, but unequally: the high edge fell from
$2.23 \times 10^{-4}$ to $5.45 \times 10^{-5}$, indistinguishable from the
interior 5--$6 \times 10^{-5}$, while the low edge improved about threefold, from
$3.20 \times 10^{-4}$ to $1.09 \times 10^{-4}$, but remains roughly twice
interior. Error then concentrated at the new low edge ($3.15 \times 10^{-4}$ for
$\omega < 0.5$, $n = 9$) but only mildly at the new high edge
($7.40 \times 10^{-5}$); mean fidelity rose to 0.999923 and the crossover moved
to $\omega = 1.53$. Edge error follows the training boundary, but
asymmetrically, which distance to the boundary alone does not explain, and the
alternative that extremal widths are intrinsically hard for a fixed receptive
field is not excluded. The ablation is one seed, its new low bin nine samples.

\section{Discussion}
\label{sec:discussion}

The spline's error is bounded not by its input's information but by
family-agnosticism: the inversion oracle recovers $\omega$ from the same 64
samples to $6.09 \times 10^{-8}$, so the information is present and a
family-blind interpolant cannot use it. That gap also bears on the
network, which reaches $1.68 \times 10^{-4}$ where inverting the same family
reaches round-off; inversion costs a fine solve per sample and is no competing
method, but bounds what the network leaves unclaimed. Its error is instead a
function of training-set design: moving a boundary moves the error, though not
symmetrically. A candidate explanation is that the relevant density is not in
$\omega$ but in the shape reproduced, since uniform sampling in $\omega$ gives
density per unit width $\sigma$ scaling as $2/\sigma^{3}$, leaving broad,
low-$\omega$ states sparsest even after widening. That is an
inference, not a tested claim; a direct test would sample uniformly in $\sigma$
or $\log\omega$.

A corollary cuts against the network: a surrogate whose accuracy tracks its
training data is weakest where this class of method is meant to be used, where
ground truth is scarcest \citep{mouli2024, gopakumar2026}; this setting isolates
that, since a physics-ordered control on identical inputs makes the departure
observable.

\section{Limitations}

\textbf{Scope.} Every sample is a harmonic ground state at fixed $x_{0} = 0$,
only $\omega$ randomized; shifted potentials, other forms (double-well and
anharmonic, implemented but untrained), and excited states are untested. Test
$\omega$ lies inside the training range, so this characterizes degradation
toward the edges, not extrapolation.

\textbf{Mechanism attribution.} The distributional explanation rests on the
low-$\omega$ asymmetry, the error shape, and one widened-range retrain. That
retrain moved edge error inward but left the new edges unequal, settling neither
the distance-to-boundary account nor the receptive-field alternative, and raising
a third, untested candidate: sampling density in shape space. It is
single-seed, its new low bin holds nine samples, and trains on 1.5 times the
samples and gradient steps of the canonical run, so its lower interior error is
not purely a range effect. The rank correlations add no independent support, and
``interior'' is unambiguous only in one dimension.

\textbf{Seed variance.} Headline numbers come from one model and one split (seed
42). Across five seeds, network mean fidelity spans 0.999832 to 0.999914 and the
crossover $\omega = 1.73$ to 2.10, so the Section~4 crossover values are one
realization and the margin varies with seed; the inverted-U with a worse low edge
appeared in every seed. The architecture was fixed a priori, not ablated.

\textbf{Fidelity is the only accuracy metric.} It is an amplitude-weighted global
overlap, insensitive to high-frequency content and boundary violation, both
visible in this model's residuals. Since the kinetic operator involves a second
derivative, a reconstruction scoring well on fidelity may score worse on energy
or $\lVert H\psi - E\psi \rVert$; no such metric was computed, and the ranking
may not survive one.

\section{Conclusion}

The two methods' errors are organized by different variables, most visibly where
the physics is easiest. Spline error is perfectly rank-ordered by $\omega$; since
inverting the coarse eigenproblem recovers $\omega$ to $6.09 \times 10^{-8}$, the
input determines the state within the family, so that failure is
representational, not informational. Network error follows an inverted-U that
five seeds reproduce, and widening the range moved error off the old edges but
not symmetrically onto the new ones, leaving training density in shape space as
the open question. Next steps: training uniform in $\sigma$ or $\log\omega$;
physics-based evaluation with energy and residual metrics; and extension to the
double-well and anharmonic potentials.

\textbf{Use of LLMs.} LLM assistance was used for drafting, analysis code
(Fourier correction, both oracles, ablation, seed replication) and auditing of
numbers and citations; all results were verified against the released code.
\textbf{Code and data.} An anonymized snapshot of the code, checkpoints, metrics
and analysis scripts accompanies this submission and regenerates every number
reported.

\begin{thebibliography}{9}

\bibitem[Raissi et al.(2019)]{raissi2019}
M.~Raissi, P.~Perdikaris, and G.~E. Karniadakis.
\newblock Physics-informed neural networks: a deep learning framework for
  solving forward and inverse problems involving nonlinear partial differential
  equations.
\newblock \emph{Journal of Computational Physics}, 378:686--707, 2019.

\bibitem[Li et al.(2021)]{li2021fno}
Z.~Li, N.~Kovachki, K.~Azizzadenesheli, B.~Liu, K.~Bhattacharya, A.~Stuart, and
  A.~Anandkumar.
\newblock Fourier neural operator for parametric partial differential equations.
\newblock In \emph{International Conference on Learning Representations (ICLR)},
  2021.

\bibitem[Dong et al.(2016)]{dong2016}
C.~Dong, C.~C. Loy, K.~He, and X.~Tang.
\newblock Image super-resolution using deep convolutional networks.
\newblock \emph{IEEE Transactions on Pattern Analysis and Machine Intelligence},
  38(2):295--307, 2016.

\bibitem[Mills et al.(2017)]{mills2017}
K.~Mills, M.~Spanner, and I.~Tamblyn.
\newblock Deep learning and the Schr\"odinger equation.
\newblock \emph{Physical Review A}, 96(4):042113, 2017.

\bibitem[Li et al.(2025)]{li2025}
C.~Li, O.~Sharir, S.~Yuan, and G.~K. Chan.
\newblock Image super-resolution inspired electron density prediction.
\newblock \emph{Nature Communications}, 16:4811, 2025.

\bibitem[Carleo and Troyer(2017)]{carleo2017}
G.~Carleo and M.~Troyer.
\newblock Solving the quantum many-body problem with artificial neural networks.
\newblock \emph{Science}, 355(6325):602--606, 2017.

\bibitem[Mouli et al.(2024)]{mouli2024}
S.~C. Mouli, D.~C. Maddix, S.~Alizadeh, G.~Gupta, A.~Stuart, M.~W. Mahoney, and
  Y.~Wang.
\newblock Using uncertainty quantification to characterize and improve
  out-of-domain learning for PDEs.
\newblock In \emph{International Conference on Machine Learning (ICML)}, 2024.

\bibitem[Gopakumar et al.(2026)]{gopakumar2026}
V.~Gopakumar, A.~Gray, J.~Oskarsson, L.~Zanisi, D.~Giles, M.~J. Kusner,
  S.~Pamela, and M.~P. Deisenroth.
\newblock Uncertainty quantification of surrogate models using conformal
  prediction.
\newblock \emph{Machine Learning: Science and Technology}, 7(1):015025, 2026.

\end{thebibliography}

\end{document}
